\begin{examplebox}
\textbf{例1：给顶角直接算面积，还要验可实现（保留原题）。}

椭圆 $x^2/25+y^2/16=1$，$\theta=60^\circ$，故
\[
 \boxed{S=16\tan30^\circ=\frac{16\sqrt3}{3}}.
\]
$c=3,b=4$，最大面积 $cb=12$，所得值小于 $12$，通过存在性初检。
进一步取
\[
 P=\left(\frac{5\sqrt{33}}9,\frac{16\sqrt3}9\right).
\]
代方程两项为 $11/27,16/27$，和为 $1$；
底乘高给 $S=3\cdot(16\sqrt3/9)=16\sqrt3/3$，
角度由本公式与半角单调性回验为 $60^\circ$。

\medskip
\textbf{例2：由面积反求顶角（保留原题）。}

椭圆 $x^2/10+y^2/5=1$ 的焦点三角形面积为 $5$，则 $b^2=5$，
\[
 \tan(\theta/2)=\frac{S}{b^2}=1
 \quad\Longrightarrow\quad\boxed{\theta=90^\circ}.
\]
回验：$c=b=\sqrt5$，最大面积恰为 $5$，故 $P=(0,\pm\sqrt5)$，
两焦半径垂直。注意所求半角必须选 $(0,90^\circ)$ 内的角。

\medskip
\textbf{例3：已知焦距、面积和角度，反求椭圆（保留原题与原 TikZ）。}

中心原点、焦点横向，焦距 $4$，$\theta=60^\circ$，面积 $4\sqrt3$。
于是 $c=2$，由 $b^2\tan30^\circ=4\sqrt3$ 得 $b^2=12$，
$a^2=b^2+c^2=16$，所以
\[
 \boxed{\frac{x^2}{16}+\frac{y^2}{12}=1}.
\]
\begin{center}\begin{tikzpicture}[font=\small,>=stealth]
\draw[black!30,rounded corners=2pt] (0,0) rectangle (14.4,8.2);
\node[anchor=west,font=\small\bfseries] at (.5,7.72) {例题的存在性：顶角 $60^\circ$ 恰好在短轴端点实现};

\begin{scope}[shift={(7.2,4.1)},scale=.6]
\draw[blue,very thick] (0,0) ellipse (4 and 3.4641016151);
\draw[orange!85!black,very thick] (-2,0)--(0,3.4641016151)--(2,0)--cycle;
\fill (-2,0) circle (3pt) node[below=6pt] {$F_1(-2,0)$};
\fill (2,0) circle (3pt) node[below=6pt] {$F_2(2,0)$};
\fill (0,3.4641016151) circle (3pt) node[above=10pt] {$P(0,2\sqrt3)$};
\node[left=6pt] at (-1,1.7320508076) {$4$};
\node[right=6pt] at (1,1.7320508076) {$4$};
\node[below=6pt] at (0,0) {$4$};
\end{scope}
\node at (7.2,1.35) {$a=4,b=2\sqrt3,c=2$；三边均为 $4$，故 $\theta=60^\circ$};

\node[anchor=west] at (.7,.55) {$S=\tfrac12\cdot4\cdot2\sqrt3=4\sqrt3=12\tan30^\circ$；两种求法一致};
\end{tikzpicture}\end{center}

取短轴端点 $P=(0,2\sqrt3)$，两焦半径和底边均为 $4$，构成等边三角形。
角度、面积与全部已知条件都能实现，另一个短轴端点同样成立。

\medskip
\textbf{例4（低门槛新题）：给点坐标，不必先求两条焦半径。}

椭圆 $x^2/25+y^2/9=1$ 上点 $P=(3,12/5)$，求焦点三角形面积及 $\tan(\theta/2)$。
先验点在椭圆上：$9/25+(144/25)/9=9/25+16/25=1$。
$a=5,b=3,c=4$，所以
\[
 \boxed{S=c|y|=4\cdot\frac{12}{5}=\frac{48}{5}},\qquad
 \boxed{\tan\frac\theta2=\frac{S}{b^2}=\frac{16}{15}}.
\]
若还要判断角型，$\tan(\theta/2)>1$，说明 $\theta>90^\circ$。
\textbf{不要先算 $PF_1,PF_2$，否则工作量反而变大。}

\medskip
\textbf{例5（易错新题）：角度公式算出了面积，但这样的点不存在。}

椭圆 $x^2/25+y^2/16=1$，有人取顶角 $\theta=90^\circ$，
便写 $S=16\tan45^\circ=16$。判断能否取到。
该椭圆 $b=4,c=3$，底乘高告诉我们\textbf{最大面积只有 $cb=12$}，
而 $16>12$。所以
\[
 \boxed{\text{此椭圆不存在顶角为 }90^\circ\text{ 的焦点三角形}}.
\]
错误不在面积公式，而在把不可能出现的 $\theta$ 当成了已存在的条件。
\end{examplebox}
\clearpage
